% Introduction to the EU – Week 9 Study Guide
% Compile: pdflatex eu-institutions-sessions.tex

\documentclass[10pt,a4paper]{article}
\usepackage[utf8]{inputenc}
\usepackage[T1]{fontenc}
\usepackage[margin=1.2cm]{geometry}
\usepackage{enumitem}
\setlist{nosep,leftmargin=*}
\usepackage{xcolor}
\usepackage{fancyhdr}
\usepackage{microtype}
\definecolor{eublue}{RGB}{0,51,153}
\pagestyle{fancy}
\fancyhf{}
\fancyfoot[L]{\small Intro EU | Week 9}
\fancyfoot[R]{\small\thepage}
\raggedright

\begin{document}
\begin{center}
{\large\bfseries\color{eublue} Introduction to the EU -- Week 9}\\[0.2em]
\footnotesize Mar 23. Housing crisis workshop + session readings. See week-6/midterm.tex \& class notes for full revision.
\end{center}\vspace{0.5em}

\section*{\color{eublue}Topics this week}
\begin{itemize}
\item Legislative / European Parliament lens: \texttt{7a Parlamentary Europe.pdf}, \texttt{7b Olsen European Parliament.pdf}, \texttt{European Parliament.pdf}
\item Housing crisis workshop: Europarl search (housing), compare group positions, and debate format mechanics
\end{itemize}

\section*{\color{eublue}Housing crisis workshop lens}
\begin{itemize}
\item Evidence to collect (for each political group):
   \begin{itemize}
   \item Latest Europarl initiatives on \emph{housing shortage / housing crisis}
   \item Lexicon used (key terms + framing words)
   \item Argument structure / proposed mechanisms (e.g., public-housing logic, affordability controls, buyer support)
   \item “Procedure logic” in the debate: when a proposal can proceed vs when it is blocked/skipped
   \end{itemize}
\item Link to the class example (Vienna housing):
   \begin{itemize}
   \item Prevention logic: avoid secondary/undesired effects
   \item Public housing requirement scale (approx. 1940s $\sim$50\%; center areas $\sim$70\% to $\sim$50\%)
   \item Effects on price + family access
   \end{itemize}
\end{itemize}

\section*{\color{eublue}Political spectrum snapshot (10th EP) -- compare stances}
\begin{itemize}
\item EPP Group: Christian democracy / conservatism; centre-right; pro-European; $184/720$ (26%)
\item S\&D Group: social democracy; centre-left; pro-European; $135/720$ (19%)
\item Renew Group: liberalism; centre to centre-right; pro-European; $76/720$ (11%)
\item Greens/EFA: green politics; centre-left to left-wing; pro-European; $53/720$ (7%)
\item The Left: democratic socialism; left-wing; soft Eurosceptic; $46/720$ (6%)
\item ECR Group: national conservatism / Atlanticism / right-wing populism; right-wing to far-right; soft Eurosceptic; $81/720$ (11%)
\item PfE (Patriots for Europe): national conservatism / right-wing populism; right-wing to far-right; Eurosceptic; $84/720$ (12%)
\item ESN: ultranationalism / ultraconservatism / right-wing populism / neo-fascism; far-right; hard Eurosceptic; $28/720$ (4%)
\item Non-Inscrits: $31/720$ (4%)
\item Vacant: $1/720$
\end{itemize}

\section*{Quick recall}
Revise: Unit 1 (four blocks), institutions. Full midterm guide in \texttt{week-6/midterm.tex}.
\end{document}
